\documentclass[10pt,a4paper]{amsart}
\usepackage[margin=19mm]{geometry}
\usepackage{amsmath,amssymb,mathtools}
\usepackage[scr=rsfs,cal=euler]{mathalfa}
\usepackage{xfrac}
\usepackage[unicode,hidelinks,bookmarks=false]{hyperref}
\newcommand{\ie}{i.e.\ }
\newcommand{\cf}{cf.\ }
\newcommand{\resp}{respectively\ }
\newcommand{\Subsection}{\subsection}
\providecommand{\nobreakdash}{\nobreak\hbox{-}}
\setlength{\emergencystretch}{2em}
\multlinegap0pt
\usepackage{bm}
\usepackage{bbm}
\usepackage{dsfont}
\usepackage{xspace}
\usepackage{cancel}
\usepackage{mathtools}
\usepackage{amsthm}
\makeatletter
\@ifundefined{lemma}{\newtheorem{lemma}{Lemma}}{}
\@ifundefined{theorem}{\newtheorem{theorem}{Theorem}}{}
\makeatother
\newcommand{\N}{\ensuremath{\mathbb{N}}\xspace}
\newcommand{\RESOCA}{RESOCA\xspace}
\newcommand{\cvres}[2]{\ensuremath{\begin{bmatrix}#1\\#2\end{bmatrix}}\xspace}
\newcommand{\cz}{\ensuremath{\underline{0}}\xspace}
\newcommand{\set}[1]{\ensuremath{\left\{#1\right\}}\xspace}
\title[The Curious Case of Reversible Elementary Second Order Cellu]{The Curious Case of Reversible Elementary Second Order Cellular Automaton 115}
\author{Enrico Formenti (Universit\'e C\^ote d'Azur, France), Supreeti Kamylia (Department of Computer Science, Engineering, Birla Institute of Technology, Mesra, Ranchi, India)}
\date{Source: arXiv:2606.13159v1 (CC BY 4.0)}
\begin{document}
\maketitle

\noindent\emph{Source and licence.} The Curious Case of Reversible Elementary Second Order Cellular Automaton 115. Version arXiv:2606.13159v1 (CC BY 4.0), distributed under the Creative Commons Attribution 4.0 International licence, as recorded in the arXiv licence metadata for this exact version and shown on the abstract page https://arxiv.org/abs/2606.13159v1. Full rights notice: licenses/arxiv-2606.13159-CC-BY-4.0.txt.\par\smallskip

\noindent\textit{Editorial scope.} Counted unit: the Theorem whose source label is \texttt{th:all-finite-cfg-are-periodic} in the article (source0.tex lines 301--304, source label \texttt{th:all-finite-cfg-are-periodic}), delivered together with the article's own complete original proof of it (source0.tex lines 498--519, one \texttt{proof} environment of 22 lines). No mathematics is written, repaired or re-derived: statement and proof are the article's own text line for line. The proof is detached from the statement in the source: the article states the result at lines 301--304 and prints its proof at lines 498--519, and the proof environment's own header names the statement, so the association is the article's own. Required context retained uncounted: lem:periodicity-1 (lines 418--437); lem:periodicity-2 (lines 438--449); lem:periodicity-3 (lines 456--475); lem:periodicity-4 (lines 477--496). Every definition, display and label of those lines is retained unchanged. Omitted from this package and not redistributed: 1--300, 305--417, 450--455, 476--476, 497--497, 520--1139. Nothing inside a retained excerpt is omitted or altered: every retained body is delivered exactly as the article's lines are written, so no omission receipt of any kind accompanies this package. Citations: the retained excerpts carry no citation command at all, so no bibliography entry of the article is transcribed and no reference is redistributed. Reference adaptation: none was needed and none was made; every reference inside the delivered unit resolves inside it, so no unresolved reference or citation is produced. Printed slips kept literal: no printed word, symbol, hypothesis or number of the counted statement or of its proof was altered. Layout and class: the article's own class is \texttt{eptcs} with packages including iftex, underscore, fontenc, breakurl, tikz, amsmath, amsthm, todonotes, cancel, subcaption, makecell, soul, refcount, xspace; the wrapper is the lane's pinned \texttt{amsart} editorial wrapper at 10pt on A4, which restates the theorem-like environments the retained text needs and adds the article's own macro definitions that the retained text needs (\texttt{\textbackslash N}, \texttt{\textbackslash RESOCA}, \texttt{\textbackslash cvres}, \texttt{\textbackslash cz}, \texttt{\textbackslash set}), with extra wrapper packages (bm, bbm, dsfont, xspace, cancel, mathtools). The delivered numbering is the unit's own. Assets: no retained span contains a figure environment, a graphics-inclusion command, a PDF-page inclusion, a table or a TikZ picture; no image file is copied into this package, the package declares no asset dependency and no image file is redistributed with it. Licence: version arXiv:2606.13159v1 of this article is distributed under the Creative Commons Attribution 4.0 International licence, as recorded in the arXiv licence metadata for this exact version and shown on the abstract page https://arxiv.org/abs/2606.13159v1. A full rights notice is delivered at licenses/arxiv-2606.13159-CC-BY-4.0.txt. Characters: every retained excerpt is pure ASCII, so the delivered engine, XeLaTeX with its default Latin Modern text font, reports no missing character.
\begin{lemma}
\label{lem:periodicity-1}
 Consider a finite configuration $b$ such that $b_i=0$ for all $i\leq0$ and another finite configuration $c$ such that $c_i=0$ for all $i<0$ and $c_0=1$. Let $f$ be the global rule of \RESOCA $115$ and $d^t=f^t\left(\cvres{b}{c}\right)_2$. Then we have for all time $t\in\N$
 \[
 \forall i\in\N\setminus\set{0,1}\quad d^t_{-i}=
 \begin{cases}
 1&\text{if\ }t\equiv1\mod3\\
 0&\text{otherwise.}
 \end{cases}
 \]
 and for all times $t\geq2$
 \[
 d^t_{-1}= \begin{cases}
 d^{t-1}_{-1}&\text{if\ }t\equiv0\mod3\\
 1&\text{if\ }t\equiv1\mod3\\
 1-(d^{t-1}_0+d^{t-2}_{-1})&\text{if\ }t\equiv2\mod3
 \end{cases}
 \]
 where addition is taken modulo $2$.
\end{lemma}

\begin{lemma}
\label{lem:periodicity-2}
 Consider two finite configurations $b$ and $c$
 such that, for both of them, all cells with an index strictly larger than $0$ are in state $0$. Let $f$ be the global rule of \RESOCA $115$ and $d^t=f^t\left(\cvres{b}{c}\right)_2$. Then we have, for all times $t\in\N$
 \[
 d^t_1=
 \begin{cases}
 1&\text{if\ }t\equiv1\mod3\\
 0&\text{otherwise.}
 \end{cases}
 \]
\end{lemma}

\begin{lemma}
\label{lem:periodicity-3}
 Consider two finite configurations $b$ and $c$ such that $b_i=c_i=0$ for all $i<0$ and $b_0=c_0=1$. Let $f$ be the global rule of \RESOCA $115$ and $d^t=f^t\left(\cvres{b}{c}\right)_2$. Then we have, for all time $t\in\N$
 \[
 \forall i\in\N\setminus\set{0}\quad d^t_{-i}=
 \begin{cases}
 1&\text{if\ }t\equiv1\mod3\\
 0&\text{otherwise.}
 \end{cases}
 \]
 and for all times $t\geq2$
 \[
 d^t_0= \begin{cases}
 d^{t-1}_0&\text{if\ }t\equiv0\mod3\\
 1&\text{if\ }t\equiv1\mod3\\
 1-(d^{t-1}_1+d^{t-2}_0)&\text{if\ }t\equiv2\mod3
 \end{cases}
 \]
 where addition is taken modulo $2$.
\end{lemma}

\begin{lemma}
\label{lem:periodicity-4}
  Consider two finite configurations $b$ and $c$ such that $b_i=c_i=0$ for all $i<0$ and $b_0=c_0=1$. Let $f$ be the global rule of \RESOCA $115$ and $d^t=f^t\left(\cvres{b}{c}\right)_2$. Then we have, for all times $t\geq3$
 \[
 \forall i\in\N\setminus\set{0,1}\quad d^t_{-i}=
 \begin{cases}
 1&\text{if\ }t\equiv1\mod3\\
 0&\text{otherwise.}
 \end{cases}
 \]
 and
 \[
 d^t_{-1}= \begin{cases}
 d^{t-1}_0&\text{if\ }t\equiv0\mod3\\
 1&\text{if\ }t\equiv1\mod3\\
 1-(d^{t-1}_1+d^{t-2}_0)&\text{if\ }t\equiv2\mod3
 \end{cases}
 \]
 where addition is taken modulo $2$.
\end{lemma}

\begin{theorem}
\label{th:all-finite-cfg-are-periodic}
    Any finite configuration is periodic for RESOCA 115.
\end{theorem}

\begin{proof}[Proof of Theorem~\ref{th:all-finite-cfg-are-periodic}]
Consider two finite configurations $b$ and $c$. If $b=c=\cz$, then it is not difficult to see that the space-time diagram of \cvres{\cz}{\cz} contains only background columns and hence has period $3$. Now assume that at least one
of $b$ or $c$ is different from \cz.
Without loss of generality we can assume that for all $i<0$, $c_i=b_i=0$
(otherwise we can choose $c'=\sigma^u(c)$ and $b'=\sigma^u(b)$, where $u=\min \cvres{b}{c}$).
By Lemma~\ref{lem:periodicity-2}, all
columns with index larger than $\max \cvres{b}{c}$ are background.
We will distinguish three disjoint cases. For each of them, we will prove
that for all $i<-1$, all columns in the space-time diagram of \cvres{b}{c} contain only background. This and the previous remark allow to deduce that
the sequence of states that we can see in the space-time diagram of initial condition \cvres{b}{c} between cell $-1$ and cell $\max\cvres{b}{c}$ is ultimately periodic. However, since RESOCA 115
is reversible, we can conclude that the sequence is periodic. Here are the
proofs for the three cases announced above:
%\begin{enumerate}
%    \item 
1) $b_0=0$ and $c_0=1$: this is Lemma~\ref{lem:periodicity-1};
%    \item 
2) $b_0=1$ and $c_0=0$: this is Lemma~\ref{lem:periodicity-4};
%    \item 
3) $b_0=1$ and $c_0=1$:
    this is Lemma~\ref{lem:periodicity-3}.
%\end{enumerate}
\end{proof}

\end{document}
